% Spectre infrarouge d'un composé inconnu de formule C4H8O : butanone (exercice 14, chapitre 8) ; spectre schématique
\begin{tikzpicture}
\begin{axis}[otIR,width=100mm,height=44mm,title={Spectre du composé inconnu}]
\addplot[otSpectre] {100 \bandeIR{2980}{25}{35} \bandeIR{2940}{20}{20} \bandeIR{1718}{18}{85} \bandeIR{1460}{15}{20}
  \bandeIR{1365}{15}{45} \bandeIR{1170}{18}{50} \bandeIR{760}{15}{20}};
\solution{
  \fill[otOrange!15] (axis cs:2850,2) rectangle (axis cs:2650,104);
  \node[pcLegende,text=otOrange!75!black,align=center] at (axis cs:2750,45) {pas de\\C--H\\d'aldéhyde};
  \node[pcLegende,text=otGreenDark,anchor=north east] at (axis cs:3010,60) {C--H};
  \node[pcLegende,text=otGreenDark,anchor=west] at (axis cs:1695,17) {C=O : cétone};}
\end{axis}
\end{tikzpicture}
